% !TeX program = pdflatex
% ALIUS Bulletin native reconstruction source.
% Compile from the ALIUS-bulletin project root. This file does not include a pre-existing article PDF.
\providecommand{\ALIUSRootPrefix}{}

\ifdefined\ALIUSIssueBuild
\else
\documentclass[a4paper,11pt]{article}
\usepackage[margin=22mm]{geometry}
\usepackage[T1]{fontenc}
\usepackage[utf8]{inputenc}
\usepackage[hidelinks]{hyperref}
\usepackage[protrusion=true,expansion=false]{microtype}
\usepackage{parskip}
\fi
\title{What is it like to be in weightlessness ?}
\author{An episode with Steven Jillings recorded by Daniel Friedman edited by Lena Coutrot, Daniel Friedman and Matthieu Koroma}
\date{ALIUS Bulletin 6}

\ifdefined\ALIUSIssueBuild
\else
\begin{document}
\fi
\maketitle
\section*{Metadata}
\textbf{Piece type:} podcast\par
\textbf{DOI:} 10.5281/zenodo.7457656\par
\textbf{Keywords:} Weightlessness, Free Fall, Spaceflight, Parabolic Flight, Vestibular System, Inversion Illusion\par
\section*{Abstract}
In this episode, Steven Jillings reviews the recent progress on the scientific study of conscious experiences in weightlessness.\par
\section*{Extracted Text}
\begingroup
\small
\noindent What is it like to be in weightlessness ? An episode with Steven Jillings Recorded by Daniel Friedman and edited by Lena Coutrot and Mahieu Koroma Cite as: Jillings, S. (2023). What is it like to be in weightlessness ? ALIUS Bullen , 6 , hps://doi.org/10.5281/zenodo.7457656 Steven Jillings Steven.Jillings@uantwerpen.be University of Antwerp, Belgium Abstract In this episode, Steven Jillings reviews the recent progress on the scienfic study of conscious experiences in weightlessness. keywords : Weightlessness, Free Fall, Spaceflight, Parabolic Flight, Vesbular System, Inversion Illusion The video and audio version of this episode are available on hps://aliusresearch.org/bullen06-jillings-podcast.html For as long as humankind can remember, we've been astounded by the mysteries of outer space. Major advancements in our understanding of the physical world in the last few centuries have led to technological developments that enable us to travel farther and farther away from Earth's surface. Although many engineering, infrastructural, and physical aspects of space travel have been successfully established, how the human body reacts and copes with such extreme conditions was and still is much less known. Of the various aspects of spaceight, weightlessness is probably the most remarkable diherence experienced by space travelers. Because biological organisms have evolved throughout a long period of time under Earth's gravitational conditions (known as 1G), their anatomical and physiological systems have developed to work optimally within this condition. Humans, therefore, experience a profound change when in weightlessness, both in terms of unconscious physiological mechanisms as well as in terms of subjective feelings. ALIUS Bulletin n6 (2023) aliusresearch.org/bulletin Steven Jillings - What is it like to be in weightlessness ? 2 Of the various aspects of spaceflight, weightlessness is probably the most remarkable difference experienced by space travelers. Starting from the rst human space mission in 1961, a new scientic research domain has emerged to unravel the ways in which the human body reacts and adjusts to weightlessness in order to ensure successful and healthy missions in space. Several approaches to studying the ehects of weightlessness on the human body exist. e most obvious way to study the ehects of weightlessness is to take measurements of astronauts while they are free-oating during their space mission. Such measurements, however, are restricted to either subjective reports or to measurements with relatively simple euipment that can be transported to space. A second approach is to perform experiments on-ground before and aer spaceight, which has the advantage of access to the full range of laboratory euipment, but has the disadvantage that measurements are not taken in a weightless condition. A common limitation for both spaceight experimental approaches is that few people go to space, leading to a slow pace of knowledge gain on the ehects of weightlessness on the human body and mind. To overcome this limitation, experimental settings on Earth have been developed to simulate various aspects of spaceight. One example is parabolic ight, which is the only way to create weightlessness on Earth. is is realized by an airplane making parabolic trajectories, where passengers experience weightlessness for 20-25 seconds during the top part of the parabola (Karmali \& Shelhamer, 2008). In this setting, researchers can conduct in-ight experiments, as well as before/aer study designs, although the passengers are only weightless for a short consecutive period. At the start of the free oating phase, disorientation and nausea may occur, which in some cases may lead to vomiting (Young et al., 1984). is is true for both spaceight and the short-term spells of weightlessness induced by parabolic ight (Glasauer \& Mittelstaedt, 1997). Weightlessness can further lead to visual and bodily illusions, such as the inversion illusion, where one's ALIUS Bulletin n6 (2023) aliusresearch.org/bulletin Steven Jillings - What is it like to be in weightlessness ? 3 sense of uprightness is lost (Gabriel et al., 1967). e reason for these experiences is that the brain is unable to make sense of conicting signals coming from the eyes and the vestibular or balance system. On Earth, our vestibular system detects head position changes, which relies in part on the existence of gravity as a vertical reference. As a conseuence, humans in weightlessness are unable to distinguish up from down based purely on their senses (Senot et al., 2012). At the start of the free floang phase, disorientaon and nausea may occur (...). Weightlessness can further lead to visual and bodily illusions, such as the inversion illusion, where one's sense of uprightness is lost. It also seems that these altered perceptions are conned to the spatial environment and not to the sense of embodiment. A study in Russian cosmonauts found that 98\% of the tested cohort experienced disorientation, while no one reported on altered perception of body ownership or embodiment (Kornilova, 1997). Additionally, the physical laws governed by gravity on Earth are hardwired within our brain to facilitate predictions of object movements in the environment and to generate uick motor responses to interact with the moving object ( Iodavina et al., 2005; Senot et al., 2012). Sensorimotor functions in a condition where the known laws of physics do not seem to apply will therefore be reuired to update and adopt new strategies, such as the trajectory of goal-directed arm movements (Gaveau et al., 2011). Finally, studies using parabolic ight have shown that weightlessness might trier altered arousal responses based on observing increased stress hormone release and elevated high-freuency neurological activity. On the other hand, aer being exposed to several parabolas, a suppression of this neurological activity was found, suesting that the participant has got somewhat familiar to the state of weightlessness (Schneider et al., 2018). Prolonged stays in outer space further trier changes in astronauts' conscious experiences. Astronauts typically have poor sleep uality due to a oating "position" during sleep (Wu et al., 2018). Accumulation of sleep ALIUS Bulletin n6 (2023) aliusresearch.org/bulletin Steven Jillings - What is it like to be in weightlessness ? 4 decit may then ahect high-demanding cognitive tasks and cause mood changes. Despite all these operational diculties and triers of the brain's autonomic responses, weightlessness is generally considered to be a profound and fun experience. Sally Ride said she believes experiencing weightlessness, along with viewing the Earth from a great distance, are for any astronaut the most fun aspects of spaceight ( https://kidadl.com/articles/inspirational-astronaut-uotes ). Experiencing weightlessness, along with viewing the Earth from a great distance, are for any astronaut the most fun aspects of spaceflight. With the prospect of future long-duration missions to the Moon and Mars, it will be important to understand how diherent G-levels (1/6G on the Moon and 1/3G on Mars) during such missions inuence conscious states and adaptability. A remarkable observation in light of transitions into diherent gravitational environments are the consistent reports that astronauts who have been to space before adapt more uickly to the space environment (Reschke et al., 1998). Upon return to Earth, these experienced astronauts also acclimatize uicker to Earth's gravity than their novice counterparts. Although there will always be a time window reuired for adaptations, a period where astronauts are most vulnerable in terms of performance, this observation suests that human beings are capable of learning to live and work in a space environment. Nevertheless, the unfamiliarity of microgravity and other gravitational states will always have a profound ehect on a vast range of brain functions, as the hardwired brain computations that have accumulated and have been optimized throughout one's life on Earth are seriously challenged during spaceight. References Gaveau, J., Paizis, C., Berret, B., Pozzo, T., Papaxanthis, C. (2011) Sensorimotor Adaptation of Point-to-Point Arm Movements aer Spaceight: e Role of Internal Representation of Gravity Force in Trajectory Planning. Journal of Neurophysiolo 106 (2): 620-29. https://doi.org/10.1152/jn.00081.2011 ALIUS Bulletin n6 (2023) aliusresearch.org/bulletin Steven Jillings - What is it like to be in weightlessness ? 5 Graybiel, A., Kello, R. S. (1967) Inversion illusion in parabolic ight: its probable dependence on otolith function. Aerospace Medical Association , 38(11):1099-103. https://doi.org/10.3357/AMHP.4195.2015 Glasauer, S. \& Mittelstaedt., H. (1997) Perception of Spatial Orientation in Diherent GLevels. Journal of Graviational Physiolo: A Journal of the International Society for Graviational Physiolo 4 (2): P5-8. Iodavina, I., Mahei, V., Bosco, G., Zago, M., Macaluso, E., Lacuaniti, F. (2005) Representation of Visual Gravitational Motion in the Human Vestibular Cortex. Science 308 (5720): 416-19. https://doi.org/10.1126/science.1107961 Karmali, F., \& Shelhamer, M. (2008). e Dynamics of Parabolic Flight: Flight Characteristics and Passenger Percepts. Aca Astronautica 63 (5-6): 594-602. https://doi.org/10.1016/j.actaastro.2008.04.009 Kornilova LN. Orientation illusions in spaceight. J Vestib Res 1997;7(6):429-39. Reschke, M. F., Bloomberg, J. J., Harm, D. L., Paloski, W. H., Layne, C., McDonald, V. (1998) Posture, Locomotion, Spatial Orientation, and Motion Sickness as a Function of Space Flight. Brain Research Reviews 28 (1-2): 102-17. https://doi.org/10.1016/S0165-0173(98)00031-9 Senot, P., Zago, M., Le Seac'h, A., Zaoui, M., Berthoz, A., Lacuaniti, F., McIntyre, J. (2012) When up Is down in 0g: How Gravity Sensing Ahects the Timing of Interceptive Actions. e Journal of Neuroscience, 32 (6): 1969-73. https://doi.org/10.1523/aNEUROSCI.3886-11.2012 Schneider, S., Brummer, V., Carnahan, H., Dubrowski, A., Askew, C. D., \& Struder, H. K. (2008). What happens to the brain in weightlessness? A rst approach by EEG tomography. Neuroimage , 42(4), 1316-1323. https://doi.org/10.1016/j.neuroimage.2008.06.010 Wu, B., Wang, Y., Wu, X., Liu, D., Xu, D., Wang, F. (2018) On-Orbit Sleep Problems of Astronauts and Countermeasures . Miliary Medical Research 5 (1): 17. https://doi.org/10.1186/s40779-018-0165-6 Young, L. R., C. M. Oman, D. G. Watt, K. E. Money, and B. K. Lichtenberg (1984) Spatial Orientation in Weightlessness and Readaptation to Earth's Gravity. Science , 225 (4658): 205-8. https://doi.org/10.1126/science.6610215 ALIUS Bulletin n6 (2023) aliusresearch.org/bulletin\par
\endgroup

\ifdefined\ALIUSIssueBuild
\clearpage
\expandafter\endinput
\fi
\end{document}
